%!TEX root = ../main.tex
% Manual title page, modelled on the one of the `tesiinfo` class used in Tesi___1
% (logo, university, department, rule, thesis type, title, advisor/candidate,
% co-advisor, rule, academic year). It replaces uaqthesis' \maketitle.
% Fill in the fields in square brackets.
%
% NOTE: the university and department names are given in their official English
% form. If the department requires the Italian wording on the title page, use
%   Università degli Studi dell'Aquila
%   Dipartimento di Ingegneria e Scienze dell'Informazione e Matematica
%   Tesi di Laurea Triennale in Ingegneria dell'Informazione
%   Relatore / Correlatore / Laureando / mat. / Anno Accademico

\begin{titlepage}
    \centering

    % --- LOGO AND UNIVERSITY ---
    \includegraphics[width=62pt]{logo_univaq.jpg}\\[1ex]
    {\bfseries\large University of L'Aquila\\[2ex]}
    {\bfseries Department of Information Engineering, Computer Science and Mathematics\\[1ex]}

    \noindent\rule[-1ex]{\textwidth}{1pt}\\[0.5ex]
    {\bfseries A thesis in partial fulfillment of the requirements for the Degree
    of Master of Science in Information Engineering\\[1ex]}

    \vfill

    % --- TITLE ---
    {\linespread{1.0}\huge\bfseries Liquid Neural Networks for\\Sensor-Agnostic Navigation Control\\}
    % \emph{\large [Subtitle]}   % uncomment and fill in if a subtitle is needed

    \vfill

    % --- ADVISOR / CANDIDATE ---
    \begin{minipage}{\textwidth}
        \raggedright
        {\bfseries Advisor: \hfill Candidate:\\[0ex]}
        Prof. Francesco Gullo \hfill {Lorenzo D'Angelo}, ID no. 302226\\[4ex]
        % --- CO-ADVISOR (delete these two lines if there is none) ---
        {\bfseries Co-Advisor:\\[0ex]}
        Rolando Innamorati\\[4ex]
    \end{minipage}

    \vfill

    % --- ACADEMIC YEAR ---
    \noindent\rule[-1ex]{\textwidth}{0.5pt}\\
    {\bfseries Academic Year 2025--26\\[1ex]}

\end{titlepage}
